\section{La membrana plasmatica: struttura}

% ---------------------------------------------------------------------------
\subsection{Aspetto e composizione}

Al microscopio elettronico la membrana appare \textbf{trilaminare} (a tre strati), dopo colorazione con osmio: l'osmio si lega preferenzialmente ai gruppi delle teste polari del doppio strato lipidico.

Componenti:
\begin{itemize}
  \item \textbf{fosfolipidi} (testa idrofilica, coda idrofobica);
  \item \textbf{colesterolo};
  \item \textbf{glicolipidi};
  \item \textbf{proteine} associate alla membrana.
\end{itemize}

% ---------------------------------------------------------------------------
\subsection{Il doppio strato fosfolipidico}

Il \textbf{fosfolipide} è una molecola anfipatica, di forma pressoché cilindrica:
\begin{itemize}
  \item \textbf{testa idrofilica}: alcol polare (es.\ colina) $+$ gruppo fosfato $+$ glicerolo;
  \item \textbf{coda idrofobica}: catene di acidi grassi.
\end{itemize}
In acqua i fosfolipidi si associano spontaneamente in un \textbf{doppio strato}: le teste idrofiliche sono rivolte verso l'acqua, le code idrofobiche verso l'interno, non a contatto con l'acqua.
\begin{itemize}
  \item confronto: i \textit{detergenti}, anfipatici ma di forma conica, in acqua si associano formando strutture sferiche (micelle);
  \item il doppio strato può ripiegarsi e chiudersi a formare una \textbf{vescicola}.
\end{itemize}

% ---------------------------------------------------------------------------
\subsection{Modelli di membrana}

\begin{itemize}
  \item \textbf{modello di Davson-Danielli} (1954): entrambe le superfici del doppio strato lipidico sono rivestite da uno strato monomolecolare di proteine, che si estendono attraverso la membrana formando canali delimitati da proteine;
  \item \textbf{modello a mosaico fluido} (Singer e Nicolson, 1972): le proteine penetrano nel doppio strato lipidico.
\end{itemize}

% ---------------------------------------------------------------------------
\subsection{Le proteine di membrana}

Tre classi:
\begin{itemize}
  \item \textbf{integrali}: attraversano il doppio strato con una o più eliche transmembrana (\textit{single pass} = una sola elica; \textit{multiple pass} = più eliche); motivi strutturali: singola $\alpha$-elica, più $\alpha$-eliche, foglietto $\beta$ ripiegato a barilotto;
  \item \textbf{periferiche}: legate con legami non covalenti ai gruppi polari delle teste lipidiche e/o a una proteina integrale;
  \item \textbf{legate ai lipidi}: unite con legami covalenti a un gruppo lipidico inserito in uno dei due foglietti (es.\ ancoraggio GPI, acido grasso, gruppo prenile).
\end{itemize}

\rubrica{Struttura}
Le regioni ad $\alpha$-elica attraversano la membrana; la proteina presenta superfici \textbf{idrofobiche} (a contatto con le code lipidiche) e \textbf{idrofiliche} (verso l'acqua o verso un canale idrofilico che attraversa la membrana). Esempi: batteriorodopsina (assorbe l'energia luminosa negli archeobatteri fotosintetici); acquaporina (proteina formata da quattro subunità che circonda un canale acquoso, avvolta da uno strato di molecole lipidiche molto aderenti).

% ---------------------------------------------------------------------------
\subsection{Il colesterolo nella membrana}

Il gruppo --OH idrofilico si inserisce nella regione polare del doppio strato, mentre la struttura ciclica (idrofobica) si estende nell'interno apolare della membrana.

% ---------------------------------------------------------------------------
\subsection{Fluidità della membrana}

\rubrica{Esperimento di Frye-Edidin}
Fusione di una cellula umana (proteine di membrana marcate con colorante fluorescente rosso) e di una cellula di topo (marcate in verde) $\rightarrow$ cellula ibrida. Dopo circa 40 minuti le proteine si mescolano sull'intera superficie: dimostra che il doppio strato di fosfolipidi è \textbf{fluido} e le proteine si muovono al suo interno.

\rubrica{Movimenti nella membrana}
La disposizione ordinata dei fosfolipidi rende la membrana un \textit{cristallo liquido}; le catene idrocarburiche sono in costante movimento e ogni molecola può spostarsi lateralmente sulla stessa faccia del doppio strato. Movimenti dei fosfolipidi:
\begin{itemize}
  \item \textbf{flessione}: $\sim 10^{-9}$ s;
  \item \textbf{spostamento laterale} (nello stesso foglietto): $\sim 10^{-6}$ s, rapido;
  \item \textbf{flip-flop} (da un foglietto all'altro): $\sim 10^{5}$ s, molto lento.
\end{itemize}

% ---------------------------------------------------------------------------
\subsection{Tecnica del freeze-fracture}

Congelamento-frattura: le cellule sono congelate rapidamente in azoto liquido e poi fratturate con un colpo secco, così da separare le due metà del doppio strato lipidico e permettere l'analisi dell'interno della membrana al microscopio elettronico. Le particelle visibili sul lato interno esposto sono \textbf{proteine integrali}. Si distinguono la \textbf{faccia E} (esterna) e la \textbf{faccia P} (protoplasmatica).

% ---------------------------------------------------------------------------
\subsection{Modello a mosaico fluido}

Schema complessivo della membrana:
\begin{itemize}
  \item doppio strato fosfolipidico (teste idrofiliche verso l'esterno, code idrofobiche all'interno);
  \item colesterolo, inserito tra i fosfolipidi;
  \item glicolipidi e glicoproteine, con i gruppi glucidici rivolti verso l'esterno della cellula;
  \item proteine integrali (tra cui le proteine canale / di trasporto) e proteine periferiche;
  \item microfilamenti del citoscheletro collegati alla membrana.
\end{itemize}
